% part: intuitionistic-logic
% chapter: introduction
% section: syntax

\documentclass[../../../include/open-logic-section]{subfiles}

\begin{document}

\olfileid[hi]{int}{int}{syn}

\olsection{अंतःप्रज्ञावादी तर्कशास्त्र की वाक्यरचना}

अंतःप्रज्ञावादी तर्कशास्त्र की वाक्यरचना प्रतिज्ञप्तिक तर्कशास्त्र जैसी ही
है। शास्त्रीय प्रतिज्ञप्तिक तर्कशास्त्र में कुछ संयोजकों को अन्य संयोजकों के
द्वारा परिभाषित किया जा सकता है; मसलन $!A \lif !B$ को
$\lnot !A \lor !B$ से, अथवा $!A \lor !B$ को
$\lnot(\lnot !A \land \lnot !B)$ से परिभाषित किया जा सकता है। अतः
शास्त्रीय तर्कशास्त्र की प्रस्तुतियाँ प्रायः कुछ संयोजकों को इन परिभाषाओं के
संक्षिप्त रूप के रूप में प्रस्तुत करती हैं। अंतःप्रज्ञावादी तर्कशास्त्र में
ऐसा नहीं है, केवल दो अपवाद हैं: $\lnot !A$ को $!A \lif \lfalse$ का संक्षिप्त
रूप परिभाषित किया जा सकता है---और प्रायः किया जाता है। तब निश्चय ही
$\lfalse$ को स्वयं परिभाषित नहीं करना चाहिए!{} साथ ही, शास्त्रीय तर्कशास्त्र
की तरह $!A \liff !B$ को $(!A \lif !B) \land (!B \lif !A)$ से परिभाषित किया
जा सकता है।

प्रतिज्ञप्तिक अंतःप्रज्ञावादी तर्कशास्त्र के !!{formula},
\emph{!!{propositional variable}ों} और प्रतिज्ञप्तिक अचर~$\lfalse$ से
\emph{तार्किक संयोजकों} का प्रयोग करके बनाए जाते हैं। हमारे पास हैं:

\begin{enumerate}
\item !!{propositional variable}ों $\Obj p_0$, $\Obj p_1$, \dots का
  !!{denumerable} समुच्चय~$\PVar$।
\item !!{falsity} का प्रतिज्ञप्तिक अचर~$\lfalse$।
\item तार्किक संयोजक: $\land$ (संयोजन), $\lor$ (वियोजन),
  $\lif$ (सशर्त)।
\item विराम-चिह्न: (, ), और अल्पविराम।
\end{enumerate}

\begin{defn}[सूत्र]
\ollabel{defn:formulas} प्रतिज्ञप्तिक अंतःप्रज्ञावादी तर्कशास्त्र के
\emph{!!{formula}ों} का समुच्चय~$\Frm[L_0]$ आगमनात्मक रूप से इस प्रकार
परिभाषित है:
\begin{enumerate}
\item $\lfalse$ एक परमाण्विक !!{formula} है।

\item प्रत्येक !!{propositional variable}~$\Obj p_i$ एक परमाण्विक !!{formula} है।

\item यदि $!A$ और $!B$ !!{formula} हैं, तो $(!A \land !B)$ एक !!{formula} है।

\item यदि $!A$ और $!B$ !!{formula} हैं, तो $(!A \lor !B)$ एक !!{formula} है।

\item यदि $!A$ और $!B$ !!{formula} हैं, तो $(!A \lif !B)$ एक !!{formula} है।

\tagitem{limitClause}{इसके अतिरिक्त कुछ भी !!{formula} नहीं है।}{}
\end{enumerate}
\end{defn}

ऊपर प्रस्तुत आदिम संयोजकों के अतिरिक्त हम निम्न \emph{परिभाषित} संकेतों का
भी उपयोग करते हैं: $\lnot$ (निषेध) और $\liff$ (!!{biconditional})। परिभाषित
संकारकों से निर्मित सूत्रों को इस प्रकार समझना है:
\begin{enumerate}
\item $\lnot !A$, $!A \lif \lfalse$ का संक्षिप्त रूप है।

\item $!A \liff !B$, $(!A \lif !B) \land (!B \lif !A)$ का संक्षिप्त
  रूप है।
\end{enumerate}

यद्यपि औपचारिक रूप से $\lnot$ को संक्षिप्त रूप माना जाता है, फिर भी कभी-कभी
हम परिभाषाओं में~$\lnot$ के लिए स्पष्ट नियम और उपवाक्य ऐसे देंगे मानो वह
आदिम हो। इसका मुख्य उद्देश्य अभ्यास-प्रश्नों को प्रस्तुत करना है।

\end{document}
